\documentclass[11pt]{article}
\usepackage[a4paper,margin=27mm]{geometry}
\usepackage[T1]{fontenc}
\usepackage{lmodern}
\usepackage{amsmath,amssymb,amsthm,booktabs}
\usepackage[colorlinks=true,linkcolor=blue,citecolor=blue,urlcolor=blue]{hyperref}
\newtheorem{theorem}{Theorem}
\newtheorem{lemma}[theorem]{Lemma}
\newtheorem{corollary}[theorem]{Corollary}
\theoremstyle{remark}
\newtheorem{remark}[theorem]{Remark}
\newcommand{\F}{\mathbb F}
\newcommand{\E}{\mathbb E}
\newcommand{\Prob}{\mathbb P}
\newcommand{\rk}{\operatorname{rank}}
\newcommand{\spn}{\operatorname{span}}
\newcommand{\ind}{\mathbf 1}
\title{Coordinatewise duality for random-access recovery times}
\author{MathIsEvenEasier\\\small Prepared with OpenAI Codex (GPT-6 Astra)}
\date{October 4, 2026}
\begin{document}
\maketitle
\begin{abstract}
For a linear code of length $n$, let $a_i$ be the expected number of independent,
uniform reads with replacement needed to recover encoded coordinate $i$.
We give a self-contained proof that the corresponding expectation for the
orthogonal dual code is $n-a_i$. In particular, recovery balance is preserved
by duality, answering Conjecture~1 of Gruica, Bar-Lev, Ravagnani and Yaakobi
under the stated sampling model. The proof also identifies $a_i/n$ with the
probability that column $i$ enters a greedy basis under a uniform random
ordering, equivalently the Shapley value of the column-rank function.
We also show that $a_i/n$ is the area under the coordinate's erasure-channel
EXIT function, relating the result to classical EXIT duality.
\end{abstract}
\noindent\textbf{Status.} A Lean certificate verifies coordinatewise duality
and recovery-balance equivalence directly for arbitrary linear codes, using
the convergent tail sum of exact uniform read-word probabilities. It also
checks that the recovery definition agrees with the generator-column span
criterion. Zero coordinates are recovered at time zero. There has been no
external peer review. The proof uses established linear algebra and
sampling identities; its connection to classical EXIT duality is explained
and attributed below.

\section{The model and the result}
Let $E=\{1,\ldots,n\}$, where $n\geq1$, and let $C\leq\F_q^n$ be a linear
code of dimension $k$, allowing $0\leq k\leq n$. Its dual is
\[
 C^\perp=\{y\in\F_q^n:\textstyle\sum_{j=1}^n x_jy_j=0
               \text{ for every }x\in C\}.
\]
Choose a generator matrix $G$ with columns $g_1,\ldots,g_n$.
For $S\subseteq E$, write $r_C(S)=\dim\spn\{g_j:j\in S\}$.
This is also the dimension of the projection of $C$ onto $S$, so it does not
depend on the generator matrix.

Let $X_1,X_2,\ldots$ be independent, uniform indices in $E$. Define
\[
 T_i(C)=\min\{t\geq0:g_i\in\spn(g_{X_1},\ldots,g_{X_t})\},
 \qquad a_i(C)=\E T_i(C).
\]
We use $\spn(\varnothing)=\{0\}$. Thus a coordinate that is identically zero
is known before any read and has recovery time zero. Equal columns retain
their distinct coordinate indices. The time to read index $i$ itself has
mean $n$ and bounds $T_i(C)$ from above, so all expectations are finite.

These are the \emph{encoded-coordinate} recovery times denoted
$\widetilde\tau_i$ in \cite{GBRY}. For a systematic generator they include the
information-coordinate times on the systematic positions. Recovery balance
means that $a_i(C)$ is independent of $i$.

\begin{theorem}[Coordinatewise duality]\label{thm:main}
For every $C\leq\F_q^n$ and every $i\in E$,
\begin{equation}\label{eq:main}
                 a_i(C)+a_i(C^\perp)=n.
\end{equation}
Consequently, $C$ is recovery balanced if and only if $C^\perp$ is recovery
balanced.
\end{theorem}

\section{A finite formula for an infinite sampling process}
For $S\subseteq E\setminus\{i\}$, set
\[
 d_C(i,S)=r_C(S\cup\{i\})-r_C(S).
\]
Adding one column increases rank by either zero or one. In particular,
$d_C(i,S)=1$ exactly when $S$ does not yet recover coordinate $i$.

\begin{lemma}[The sampling identity]\label{lem:sampling}
For every coordinate $i$,
\begin{equation}\label{eq:sampling}
 a_i(C)=\sum_{S\subseteq E\setminus\{i\}}
              \frac{d_C(i,S)}{\binom{n-1}{|S|}}.
\end{equation}
\end{lemma}
\begin{proof}
Continue drawing indices even after coordinate $i$ is recovered.
Let $D_s$ be the time at which $s$ distinct indices have first been seen,
with $D_0=0$, and let $S_s$ be this set of indices. For $0\leq s<n$, put
$W_s=D_{s+1}-D_s$. All $D_s$ are finite almost surely.

Between $D_s$ and the next new index the span does not change. Therefore,
pathwise,
\begin{equation}\label{eq:path}
 T_i(C)=\sum_{s=0}^{n-1}
       \ind_{\{g_i\notin\spn(g_j:j\in S_s)\}}W_s.
\end{equation}
Conditional on the entire history through $D_s$, the waiting time $W_s$ is
geometric with success probability $(n-s)/n$, hence conditional mean
$n/(n-s)$. This conditional-mean statement is sufficient; independence
between the indicator and the waiting time is not being assumed.

By symmetry $S_s$ is uniform among all $s$-subsets of $E$.
Let $b_i(s)$ count those $s$-subsets that do not recover $i$.
Such subsets cannot contain $i$. Taking expectations in \eqref{eq:path}
gives
\[
 a_i(C)=\sum_{s=0}^{n-1}\frac{n}{n-s}
                    \frac{b_i(s)}{\binom ns}
        =\sum_{s=0}^{n-1}\frac{b_i(s)}{\binom{n-1}s}.
\]
The last equality uses
$(n-s)\binom ns=n\binom{n-1}s$. Expanding $b_i(s)$ as the sum of the rank
increments proves \eqref{eq:sampling}. The factor $n/(n-s)$ is what accounts
for repeated reads.
\end{proof}

\section{Complementary recovery sets in the dual}
\begin{lemma}[Dual ranks]\label{lem:rank}
For every $A\subseteq E$,
\begin{equation}\label{eq:rank}
 r_{C^\perp}(A)=|A|-k+r_C(E\setminus A).
\end{equation}
If $S\subseteq E\setminus\{i\}$ and
$U=E\setminus(S\cup\{i\})$, then
\begin{equation}\label{eq:increment}
                   d_{C^\perp}(i,U)=1-d_C(i,S).
\end{equation}
\end{lemma}
\begin{proof}
For any $B\subseteq E$, the vectors in $C^\perp$ supported on $B$ correspond
to the orthogonal complement of the projection of $C$ onto $B$, inside
$\F_q^B$. Their dimension is $|B|-r_C(B)$.
The kernel of projection from $C^\perp$ onto $A$ consists exactly of the
vectors supported on $E\setminus A$. Rank-nullity, and
$\dim C^\perp=n-k$, therefore give
\[
 r_{C^\perp}(A)=(n-k)
                 -\bigl(n-|A|-r_C(E\setminus A)\bigr),
\]
which proves \eqref{eq:rank}. Apply it at $U\cup\{i\}$ and $U$ and subtract:
\begin{align*}
 d_{C^\perp}(i,U)
 &=\bigl(|U|+1-k+r_C(S)\bigr)
       -\bigl(|U|-k+r_C(S\cup\{i\})\bigr)\\
 &=1-d_C(i,S).
\end{align*}
\end{proof}

\begin{proof}[Proof of Theorem~\ref{thm:main}]
The map $S\mapsto E\setminus(S\cup\{i\})$ is a bijection of the subsets
of $E\setminus\{i\}$, sending size $s$ to $n-1-s$. The denominator in
\eqref{eq:sampling} is unchanged under this map. Consequently,
\begin{align*}
 a_i(C)+a_i(C^\perp)
 &=\sum_{S\subseteq E\setminus\{i\}}
       \frac{d_C(i,S)+d_{C^\perp}(i,E\setminus(S\cup\{i\}))}
            {\binom{n-1}{|S|}}\\
 &=\sum_{s=0}^{n-1}\frac{\binom{n-1}s}{\binom{n-1}s}=n.
\end{align*}
Thus $a_i(C^\perp)=n-a_i(C)$ for every $i$. One profile is constant exactly
when the other is constant, which is Conjecture~1 of \cite{GBRY}.
\end{proof}

\section{The connection with classical EXIT duality}\label{sec:exit}
Fix $i$ and retain each coordinate other than $i$, independently, with
probability $1-\varepsilon$. Let $S_\varepsilon$ be the retained set and put
\begin{equation}\label{eq:exit}
 h_{C,i}(\varepsilon)
 =\Prob\{S_\varepsilon\text{ does not recover }i\}
 =\sum_{S\subseteq E\setminus\{i\}}
 d_C(i,S)(1-\varepsilon)^{|S|}\varepsilon^{n-1-|S|}.
\end{equation}
This polynomial is defined on $[0,1]$, including its endpoints. The target
coordinate itself is always withheld in this experiment.

For a uniform random codeword $Z\in C$, send the other coordinates through
independent erasure channels with erasure probability $\varepsilon$ and
call the output $Y_{-i}$. With entropy measured in base $q$,
\[
                 h_{C,i}(\varepsilon)=H_q(Z_i\mid Y_{-i}).
\]
Indeed, after any fixed observation the possible codewords form an affine
coset of the subspace vanishing on the observed coordinates. Its projection
onto coordinate $i$ is either a singleton or all of $\F_q$, uniformly
distributed. The corresponding conditional entropy is respectively zero
or one. This also covers identically zero coordinates. For binary codes
this is the usual entropy convention for the coordinate EXIT function;
mutual-information conventions require the corresponding change of variables.

\begin{corollary}[Recovery time as an EXIT area]\label{cor:exit}
For every coordinate,
\begin{equation}\label{eq:exit-area}
             \frac{a_i(C)}{n}=\int_0^1 h_{C,i}(\varepsilon)\,d\varepsilon.
\end{equation}
Moreover,
\begin{equation}\label{eq:exit-dual}
             h_{C,i}(\varepsilon)+h_{C^\perp,i}(1-\varepsilon)=1.
\end{equation}
Thus integrating the second identity proves Theorem~\ref{thm:main}.
\end{corollary}
\begin{proof}
For $0\leq s\leq n-1$, elementary beta integration gives
\[
 \int_0^1(1-\varepsilon)^s\varepsilon^{n-1-s}\,d\varepsilon
 =\frac{s!(n-1-s)!}{n!}=\frac{1}{n\binom{n-1}s}.
\]
Integrate the finite sum in \eqref{eq:exit} and apply
Lemma~\ref{lem:sampling}. For the second identity, pair each observed set
$S$ with $U=E\setminus(S\cup\{i\})$. At erasure parameter
$1-\varepsilon$, the probability of observing $U$ equals the probability
of observing $S$ at parameter $\varepsilon$. Their failure indicators add
to one by \eqref{eq:increment}. The probabilities of all $S$ sum to one,
proving \eqref{eq:exit-dual}. A change of variable in the dual integral
then gives $a_i(C)/n+a_i(C^\perp)/n=1$.
\end{proof}

\paragraph{Attribution.}
Erasure-channel EXIT duality predates this note: Ashikhmin, Kramer and
ten Brink prove a duality property in \cite[Section IV-E, Theorem 4]{AKB}.
Renes gives a broader channel-duality framework in
\cite[Section 4.5, Theorem 4 and its coordinatewise proof]{Renes}.
Equation~\eqref{eq:exit-dual} is this classical phenomenon in our failure
probability notation, and is not claimed as a new duality theorem.
Equation~\eqref{eq:exit-area} explicitly connects it to the uniform-read
model of the recovery-balance conjecture. Together with
\eqref{eq:greedy}, it identifies three quantities:
\[
 \frac{a_i(C)}n
   =\int_0^1 h_{C,i}(\varepsilon)\,d\varepsilon
   =\phi_i(r_C).
\]
In particular, recovery balance means equality of the coordinate EXIT
areas. It does not require equality of the entire EXIT functions.
This section is a written derivation; EXIT functions and integration have
not been added to the Lean certificate.

\section{A permutation interpretation and consequences}
Choose a uniformly random permutation $\pi$ of $E$. Process its columns in
order, retaining a column precisely when it increases the rank of the
retained columns. Let $B_\pi(C)$ be the resulting basis and let $P_i(\pi)$
be the indices preceding $i$ in the permutation.

\begin{corollary}[Greedy inclusion probability]\label{cor:greedy}
For every $i$,
\begin{equation}\label{eq:greedy}
 a_i(C)=n\Prob\{i\in B_\pi(C)\}
       =n\,\phi_i(r_C),
\end{equation}
where $\phi_i(r_C)$ is the Shapley value of the rank set function.
In particular,
\begin{equation}\label{eq:sum}
                     \sum_{i=1}^n a_i(C)=nk.
\end{equation}
\end{corollary}
\begin{proof}
For each fixed $S\subseteq E\setminus\{i\}$ there are
$|S|!(n-1-|S|)!$ permutations with $P_i(\pi)=S$. Hence
\[
 \Prob\{P_i(\pi)=S\}=\frac1{n\binom{n-1}{|S|}}.
\]
The retained columns span all previously processed columns, so $i$ is
retained exactly when $d_C(i,P_i(\pi))=1$. Equation~\eqref{eq:sampling}
proves the first equality in \eqref{eq:greedy}. The second is the usual
average-marginal-contribution definition of the Shapley value \cite{Shapley}.
Every $B_\pi(C)$ has exactly $k$ elements. Summing its inclusion
probabilities proves \eqref{eq:sum}, also obtained in \cite{GBRY}.
\end{proof}

This gives a finite way to compute or estimate an entire recovery profile:
average the greedy inclusion indicators and multiply by $n$.
Each permutation requires at most $n$ rank-update decisions and has no
unbounded waiting for repeated reads. The induced distribution on bases
need not be uniform. With independent sampled permutations this gives an
unbiased estimator; no empirical speed or accuracy benchmark is claimed here.
There is also an exact pairing:
\[
 B_{\operatorname{reverse}(\pi)}(C^\perp)=E\setminus B_\pi(C).
\]
Indeed, the predecessor set of $i$ in the reversed order is the complement
of $P_i(\pi)\cup\{i\}$, so \eqref{eq:increment} makes the two inclusion
indicators complementary.

\begin{corollary}[Recovery profiles]\label{cor:profiles}
The mean of the coordinate expectations is $k$ for $C$ and $n-k$ for
$C^\perp$, and
\[
 a_i(C^\perp)-(n-k)=-\bigl(a_i(C)-k\bigr).
\]
In particular,
\[
 \max_i a_i(C^\perp)=n-\min_i a_i(C),\qquad
 a_i(C)<a_j(C)\iff a_i(C^\perp)>a_j(C^\perp).
\]
The two profiles have the same range and the same variance across
coordinates. If $C=C^\perp$ on the same coordinate set, every coordinate
has expectation $n/2$.
\end{corollary}
\begin{proof}
Use \eqref{eq:main} and \eqref{eq:sum}. Centering subtracts $k$ and $n-k$,
respectively, and negates each deviation; squares and absolute values are
preserved. If $C=C^\perp$, \eqref{eq:main} becomes $2a_i(C)=n$.
\end{proof}

\section{An unbalanced example, checked directly}
Over $\F_2$, consider the generator matrices
\[
 G=\begin{pmatrix}1&1&1&0&0\\0&0&0&1&1\end{pmatrix},\qquad
 H=\begin{pmatrix}1&1&0&0&0\\1&0&1&0&0\\0&0&0&1&1\end{pmatrix}.
\]
Their ranks are $2$ and $3$, and $GH^{\mathsf T}=0$, so their row spaces
are dual codes of length five. All columns of both matrices are nonzero.
\begin{center}
\begin{tabular}{lccc}
\toprule
Target coordinates & $a_i(C)$ & $a_i(C^\perp)$ & Sum\\
\midrule
$1,2,3$ & $5/3$ & $10/3$ & $5$\\
$4,5$ & $5/2$ & $5/2$ & $5$\\
\bottomrule
\end{tabular}
\end{center}
For $G$, recovering any of the first three coordinates requires hitting
their group, a geometric waiting time of mean $5/3$. The other group gives
mean $5/2$.

For $H$, target $1$ is recovered by reading index $1$, or by reading both
indices $2$ and $3$. Indices $4,5$ do not help with it. Inclusion-exclusion
therefore gives, for every integer $t\geq0$,
\[
 \Prob\{T_1(C^\perp)>t\}=2(3/5)^t-(2/5)^t.
\]
Summing this tail yields $2/(1-3/5)-1/(1-2/5)=10/3$.
The same calculation applies to targets $2,3$. Targets $4,5$ are recovered
when either member of their group appears, giving $5/2$.
Thus the table is checked without invoking Theorem~\ref{thm:main}.
The easiest group for $C$ is the hardest group for $C^\perp$.
The EXIT description gives the same check: for $i=1,2,3$,
\[
 h_{C,i}(\varepsilon)=\varepsilon^2,\qquad
 h_{C^\perp,i}(\varepsilon)=2\varepsilon-\varepsilon^2.
\]
The areas are $1/3$ and $2/3$; multiplying by five yields the first row of
the table. For $i=4,5$ both EXIT functions equal $\varepsilon$, of area
$1/2$.

\section{Scope and boundary cases}
\paragraph{Zero coordinates.}
The sampling lemma includes $S=\varnothing$. For a zero column every rank
increment is zero, so its expectation is zero. The dual coordinate is a
coloop: it can only be recovered by reading its own index, giving mean $n$.
This handles all degenerate cases in Theorem~\ref{thm:main}, including the
zero code and the full-space code. The displayed formula (4) in
\cite{GBRY}, with its sum starting at $s=1$, requires an empty-set correction
under our time-zero convention for a zero target. The source does not
explicitly settle this boundary convention; our proof states it and derives
the required formula directly.

If one instead insists on taking at least one read even for a known zero,
the coordinatewise equation needs that convention's correction, but
recovery balance is still preserved by duality. Here are the boundary cases.
A proper nonzero code with a zero column has expectation $1$ there and
expectation greater than $1$ at any nonzero column: drawing the zero column
first prevents immediate recovery. Hence it is unbalanced. Its dual has a
coloop of mean $n$ and, since it is not the full-space code, a noncoloop.
A nonzero noncoloop has an alternative recovery set, which is read before
the target with positive probability, so its expectation is strictly below
$n$; a zero column has expectation $1<n$ here. Thus the dual is unbalanced
as well. Interchange the two codes if only the dual has a zero column.
The zero and full-space codes are separately balanced, including $n=1$.
When neither code has zero columns, the conventions agree.

\paragraph{Uniform reads are essential.}
Take the binary repetition code $C=\spn\{(1,1,1)\}$ and its parity-check
dual, and sample indices with probabilities $(1/2,1/4,1/4)$. Any read
recovers target $1$ in $C$, so its mean is $1$. In $C^\perp$, target $1$
is recovered by reading index $1$ or both indices $2,3$. For $t\geq1$,
failure means that all $t$ reads were index $2$, or all were index $3$.
Therefore its mean is
\[
 1+\sum_{t=1}^\infty 2(1/4)^t=5/3,
 \qquad 1+5/3=8/3\ne3.
\]
The theorem consequently cannot be applied unchanged to nonuniform
sequencing. It also makes no assertion about read errors, simultaneous
recovery of all coordinates, or decoding runtime.

\paragraph{Matroid form.}
The same proof applies to every finite matroid, with recovery defined by
membership in the closure of the drawn elements. Lemma~\ref{lem:sampling}
uses only closure and uniform sampling, and \eqref{eq:increment} follows
from the standard dual-rank formula. Representability over a finite field
is needed only to interpret the result as a statement about linear codes.

\paragraph{Practical use.}
Within the uniform error-free model, the theorem transfers a computed
recovery profile to the dual and reverses the ordering of easy and difficult
coordinates. The greedy interpretation supplies a bounded-work unbiased
estimator per permutation. Neither statement by itself designs a better
physical DNA storage system.

\section{Relation to earlier work}
Notation~1 and Definition~2 of \cite{GBRY} specify the encoded-coordinate
sampling model and recovery balance used here. Theorem~\ref{thm:main}
proves their Conjecture~1 through a coordinatewise identity. Section~\ref{sec:exit}
explains how classical erasure-channel EXIT duality yields the same result
once its relation to expected read counts is established.

The duality in \cite[Section 4]{BRY} expresses the time for full recovery
using subcode data of the dual. Its observable is different from the
single-coordinate time considered here. Two further related works were
screened: \cite{Bodur} on expectations and distributions for random access,
and \cite{WY}, including its August 2026 revision, on algorithms,
constructions and bounds.

The objective of this note is to resolve the stated conjecture with a
checkable proof using available mathematics. The earlier results provide
the tools and context for that resolution. The accompanying Lean certificate
checks the code theorem in the precise model described next. Literature
search records are preserved in \texttt{publication-review.json}; establishing
publication priority is not an objective or a condition for completion.

\section{What Lean has checked}
The code theorem quantifies over an arbitrary field $K$, a finite nonempty
coordinate type $E$, and an arbitrary subspace $C\leq K^E$. Thus it applies
to every finite field and every code length, without restricting the code
family or enumerating examples. Its formal orthogonal dual satisfies
\[
 y\in C^\perp\quad\Longleftrightarrow\quad
           \sum_{j\in E}x_jy_j=0\quad\text{for every }x\in C.
\]

\paragraph{Recoverability and generator columns.}
The formal recovery predicate is
\[
 R_C(S,i)\quad\Longleftrightarrow\quad
 \bigl(\forall x\in C,\ x|_S=0\ \Longrightarrow\ x_i=0\bigr).
\]
Lean proves that this is equivalent to any two codewords agreeing on $S$
also agreeing at $i$. It further proves equivalence with the existence of
coefficients $\lambda_j$ such that
\[
                   x_i=\sum_{j\in S}\lambda_j x_j
                   \qquad\text{for every }x\in C.
\]
For a code generated by columns $g_j$, \texttt{CodeGenerator.lean} proves
$R_C(S,i)\Longleftrightarrow g_i\in\spn\{g_j:j\in S\}$, exactly the column
criterion in the sampling model. No full-row-rank assumption is needed for
this equivalence.

\paragraph{A direct proof of complementary recovery.}
The formal code proof establishes the following equivalence directly,
without assuming a code-to-matroid representation theorem. For
$S\subseteq E\setminus\{i\}$ and $U=E\setminus(S\cup\{i\})$,
\[
                        R_{C^\perp}(U,i)\Longleftrightarrow\neg R_C(S,i).
\]
If $R_C(S,i)$ holds, use its linear decoder to construct
\[
                         y=e_i-\sum_{j\in S}\lambda_j e_j.
\]
Then $y\in C^\perp$, $y|_U=0$, and $y_i=1$. This witnesses failure of
$R_{C^\perp}(U,i)$. Conversely, failure of $R_C(S,i)$ gives $x\in C$ with
$x|_S=0$ and $x_i\ne0$. For every $y\in C^\perp$ with $y|_U=0$,
orthogonality reduces to $x_iy_i=0$, hence $y_i=0$. This proves
$R_{C^\perp}(U,i)$. The existence of the linear decoder, and both directions
of this argument, are checked in \texttt{CodeLinear.lean}.

\paragraph{Exact sampling and the final theorem.}
Put $n=|E|$. For each $t\geq0$, Lean identifies the formal failure
probability with the exact fraction
\[
 q_{C,i}(t)=\frac{\bigl|\{w\in E^t:
           \neg R_C(\{w_1,\ldots,w_t\},i)\}\bigr|}{n^t}.
\]
Words are ordered and allow repeated indices; all $n^t$ words receive the
same weight. The formal expected read count is defined by
$A_i(C)=\sum_{t=0}^{\infty}q_{C,i}(t)$. Lean proves this series converges,
proves the finite subset formula, and then proves
\[
 A_i(C)+A_i(C^\perp)=n,
 \qquad
 \bigl(A_i(C)\text{ is constant in }i\bigr)
 \Longleftrightarrow
 \bigl(A_i(C^\perp)\text{ is constant in }i\bigr).
\]
The main entry points in \texttt{CodeResult.lean} are
\texttt{codeExpectedReads\_dual} and
\texttt{codeRecoveryBalanced\_dual\_iff}.
\texttt{CodeProbability.lean} checks the failure-count interpretation.
The earlier sampling modules prove convergence and the subset formula,
and \texttt{Pairing.lean} pairs complementary subsets.
The separate general finite-matroid theorem remains available in
\texttt{Result.lean}.

\paragraph{Verification and trust.}
All thirteen modules, including both axiom audits, compiled with warnings
treated as errors. The certificate uses Lean 4.34.0 and mathlib revision
\texttt{5ed2965256430c3649e86755f9576b54eca72435}.
The axiom audit reports only \texttt{propext}, \texttt{Classical.choice},
and \texttt{Quot.sound}. The proof uses no custom axioms, admitted goals,
\texttt{unsafe} declarations or \texttt{native\_decide}. A deliberately
ill-typed proof of \texttt{False} was rejected in a separate control run.
This is a standard Lean kernel check with pinned upstream mathlib cache
artifacts, not an independent implementation of the kernel. Substantial
compilations run only in Azure, with a 12 GiB process-group memory cap,
a maximum 35-minute worker lifetime, 90-second module timeouts and an
independent cloud deletion guard. The code-certificate job is
\texttt{83cb03dc071a4251}; its build, axiom audit, negative control, source
hashes and cleanup evidence are recorded in
\texttt{formal/verification-code-azure.json}.
The complete result archive was retrieved and all fifteen source files
(including the two negative controls) matched the local sources and the
final build hashes. Deletion of both job resource groups was confirmed at
08:16 UTC on October 4, 2026.
The earlier matroid certificate is preserved separately in
\texttt{formal/verification-azure.json}.

\paragraph{Exact scope of the probability formalization.}
The certificate uses the explicit convergent finite-word tail sum above,
which is the standard expression $\E T=\sum_{t\geq0}\Prob(T>t)$ for a
nonnegative integer-valued recovery time. It does not separately construct
an infinite product probability space and a hitting-time random variable
on that space. The written sampling argument supplies the usual stochastic
interpretation. The Shapley and EXIT interpretations, worked numerical examples and
alternative convention requiring a read even for a known zero remain
outside the present Lean certificate. The linear-code duality and
recovery-to-column-span bridges are now included in the certificate.

\begin{thebibliography}{9}
\bibitem{GBRY}
A. Gruica, D. Bar-Lev, A. Ravagnani and E. Yaakobi,
\emph{A Combinatorial Perspective on Random Access Efficiency for DNA Storage},
arXiv:2401.15722v3, September 24, 2025.
\href{https://arxiv.org/html/2401.15722v3}{Full text}.
\href{https://doi.org/10.1109/TIT.2025.3623202}{IEEE DOI}.
The target is Conjecture~1 in Section~5.3.

\bibitem{Shapley}
L. S. Shapley, \emph{A Value for $n$-Person Games},
in \emph{Contributions to the Theory of Games II},
Annals of Mathematics Studies 28, Princeton University Press, 1953,
pp.~307--317. Earlier version: RAND P-295, 1952,
\url{https://www.rand.org/pubs/papers/P295.html}.

\bibitem{AKB}
A. Ashikhmin, G. Kramer and S. ten Brink,
\emph{Extrinsic Information Transfer Functions: Model and Erasure Channel
Properties}, IEEE Transactions on Information Theory 50(11), 2657--2673,
2004. \href{https://doi.org/10.1109/TIT.2004.836693}{DOI}.

\bibitem{Renes}
J. M. Renes, \emph{Duality of channels and codes}, IEEE Transactions on
Information Theory 64, 577, 2018. arXiv:1701.05583v2, July 10, 2017.
\href{https://arxiv.org/html/1701.05583v2}{Full text}.
\href{https://doi.org/10.1109/TIT.2017.2754921}{DOI}.

\bibitem{BRY}
M. Bertuzzo, A. Ravagnani and E. Yaakobi,
\emph{The DNA Coverage Depth Problem: Duality, Weight Distributions, and
Applications}, arXiv:2603.06489v1, March 6, 2026.
\url{https://arxiv.org/html/2603.06489v1}.
This related duality result concerns full recovery.

\bibitem{Bodur}
\c{S}. Bodur, S. Lia, H. H. L\'opez, R. Ludhani, A. Ravagnani and
L. Seccia, \emph{The Random Variables of the DNA Coverage Depth Problem},
arXiv:2507.20645v1, July 28, 2025.
\url{https://arxiv.org/html/2507.20645v1}.

\bibitem{WY}
C. Wang and E. Yaakobi, \emph{Random Access in DNA Storage: Algorithms,
Constructions, and Bounds}, arXiv:2601.07053v2, August 12, 2026.
\url{https://arxiv.org/html/2601.07053v2}.
\end{thebibliography}
\end{document}
